% Modular TeX snippet for the 253b final paper.
% Purpose: restore detailed calculations that were compressed in the first modular pass.

\providecommand{\dd}{\mathrm{d}}
\providecommand{\Tr}{\operatorname{Tr}}
\providecommand{\tr}{\operatorname{tr}}
\providecommand{\CS}{\operatorname{CS}}
\providecommand{\Hom}{\operatorname{Hom}}
\providecommand{\U}{\operatorname{U}}
\providecommand{\SU}{\operatorname{SU}}
\providecommand{\Lk}{\operatorname{Lk}}
\providecommand{\SLk}{\operatorname{SLk}}
\providecommand{\Lie}{\operatorname{Lie}}
\providecommand{\id}{\operatorname{id}}

\section{Worked Derivations}
\label{sec:dense-derivation-expansion}

The preceding sections state the structure of Chern--Simons theory in a compressed form. This section deliberately slows down. Its purpose is to keep the paper from becoming a list of elegant slogans. We spell out the calculations that make the slogans true: how differential forms encode ordinary component formulae, how the Wess--Zumino integral becomes an integer, how the canonical torus Hilbert space is counted, how abelian Wilson loops turn into linking numbers, and how the $K$-matrix response formulas follow from integrating out internal gauge fields.

The intended rhythm is: say what is being computed, do the algebra without skipping convention-sensitive signs, and then say what the result means physically. If a derivation also appears in the synthesis module, the version here is the worked version: it is where a reader should go when they want to know exactly where a factor of $2$, a boundary term, or a normalization has entered.

% FIGURE FLAG [F-worked-map]: a dependency map from the worked derivations to level quantization, torus Hilbert spaces, Wilson phases, framing, K-matrix response, and edge modes would make this module easier to navigate.

\subsection{Component Dictionary}
\label{subsec:dense-component-ff}

Let
\begin{equation}
  F=\frac12F_{\mu\nu}\,\dd x^\mu\wedge \dd x^\nu.
  \label{eq:dense-F-components}
\end{equation}
Then
\begin{align}
  F\wedge F
  &=
  \frac14
  F_{\mu\nu}F_{\rho\sigma}
  \,\dd x^\mu\wedge\dd x^\nu\wedge\dd x^\rho\wedge\dd x^\sigma
  \notag\\
  &=
  \frac14
  \epsilon^{\mu\nu\rho\sigma}
  F_{\mu\nu}F_{\rho\sigma}
  \,\dd^4x.
  \label{eq:dense-ff-component}
\end{align}
After tracing over gauge indices,
\begin{equation}
  \tr(F\wedge F)
  =
  \frac14
  \epsilon^{\mu\nu\rho\sigma}
  \tr(F_{\mu\nu}F_{\rho\sigma})\,\dd^4x.
  \label{eq:dense-trff-component}
\end{equation}
This is the bridge to the component expression in anomaly calculations. If
\begin{equation}
  \widetilde F^{\mu\nu}
  =
  \frac12\epsilon^{\mu\nu\rho\sigma}F_{\rho\sigma},
  \label{eq:dense-dual-F}
\end{equation}
then
\begin{equation}
  \tr(F_{\mu\nu}\widetilde F^{\mu\nu})
  =
  \frac12\epsilon^{\mu\nu\rho\sigma}
  \tr(F_{\mu\nu}F_{\rho\sigma}).
  \label{eq:dense-F-Ftilde}
\end{equation}
Comparing \eqref{eq:dense-trff-component} and \eqref{eq:dense-F-Ftilde}, the form density differs from the usual $F\widetilde F$ component density by a factor of $1/2$:
\begin{equation}
  \tr(F\wedge F)
  =
  \frac12
  \tr(F_{\mu\nu}\widetilde F^{\mu\nu})\,\dd^4x.
  \label{eq:dense-ff-ftilde-relation}
\end{equation}
The exact numerical coefficient depends on whether the component convention includes $\tr(T_aT_b)=\delta_{ab}$ or $\frac12\delta_{ab}$, and on whether the gauge coupling is included in $A$. The important invariant statement is independent of these bookkeeping choices: the four-form $\tr(F\wedge F)$ is closed and locally exact, and its local primitive is the Chern--Simons three-form.

This dictionary matters because the anomaly literature and the topological field theory literature often use different languages for the same object. In components one sees $F_{\mu\nu}\widetilde F^{\mu\nu}$; in forms one sees $\tr(F\wedge F)$. Once the factor in \eqref{eq:dense-ff-ftilde-relation} is fixed, the divergence of the Chern--Simons current and the identity $\dd\omega_{\CS}=\tr(F\wedge F)$ become the same statement.

\subsection{Quartic Trace}
\label{subsec:dense-quartic-zero}

The cancellation of the quartic term in
\begin{equation}
  \tr\left((\dd A+A\wedge A)\wedge(\dd A+A\wedge A)\right)
  \label{eq:dense-ff-quartic-context}
\end{equation}
is a useful example of how trace cyclicity and graded signs cooperate. Let
\begin{equation}
  Q=\tr(A\wedge A\wedge A\wedge A).
  \label{eq:dense-quartic-Q}
\end{equation}
Cyclicity of the trace lets us move the first matrix factor to the end without changing the trace as a matrix trace. But the form part is made of one-forms, so moving that first $A$ past three other one-forms gives $(-1)^3=-1$. Therefore
\begin{equation}
  Q
  =
  -Q.
  \label{eq:dense-quartic-sign}
\end{equation}
Over a field of characteristic zero this implies $Q=0$. This small algebraic fact is what makes
\begin{equation}
  \dd\,\tr\left(A\wedge\dd A+\frac23A\wedge A\wedge A\right)
  =
  \tr(F\wedge F)
  \label{eq:dense-transgression-recalled}
\end{equation}
work with exactly the coefficient $2/3$.

It is worth emphasizing what is not being claimed. The form $A\wedge A\wedge A\wedge A$ is not zero as a matrix-valued form in general. What vanishes is its trace. This distinction matters because the nonabelian product contains both antisymmetric form data and noncommuting matrix data.

This is a good example of why a terse derivation can mislead. If one forgets the trace, the statement looks false; if one forgets the graded sign, it looks mysterious. The cancellation uses both ingredients at once: matrix cyclicity and one-form antisymmetry.

\subsection{Action Variation}
\label{subsec:dense-full-variation}

Start from
\begin{equation}
  \omega_{\CS}(A)
  =
  \tr\left(A\wedge\dd A+\frac23A\wedge A\wedge A\right).
  \label{eq:dense-cs-form}
\end{equation}
Under $A\mapsto A+\delta A$,
\begin{align}
  \delta\omega_{\CS}
  &=
  \tr\left(
  \delta A\wedge\dd A
  +
  A\wedge\dd(\delta A)
  +
  \frac23\delta(A\wedge A\wedge A)
  \right).
  \label{eq:dense-var-start}
\end{align}
The cubic variation is
\begin{align}
  \delta(A\wedge A\wedge A)
  &=
  \delta A\wedge A\wedge A
  +
  A\wedge\delta A\wedge A
  +
  A\wedge A\wedge\delta A.
  \label{eq:dense-cubic-var}
\end{align}
Inside the trace, each term equals $\tr(\delta A\wedge A\wedge A)$. Moving the one-form $\delta A$ through two one-forms costs no sign, and cyclicity handles the matrix order. Hence
\begin{equation}
  \tr\,\delta(A\wedge A\wedge A)
  =
  3\tr(\delta A\wedge A\wedge A).
  \label{eq:dense-cubic-var-trace}
\end{equation}
The middle term in \eqref{eq:dense-var-start} is integrated by parts using
\begin{equation}
  \dd\,\tr(A\wedge\delta A)
  =
  \tr(\dd A\wedge\delta A-A\wedge\dd(\delta A)).
  \label{eq:dense-var-boundary-identity}
\end{equation}
Solving for $\tr(A\wedge\dd(\delta A))$ gives
\begin{equation}
  \tr(A\wedge\dd(\delta A))
  =
  \tr(\dd A\wedge\delta A)
  -
  \dd\,\tr(A\wedge\delta A).
  \label{eq:dense-var-middle}
\end{equation}
Since $\dd A$ is a two-form and $\delta A$ is a one-form,
\begin{equation}
  \tr(\dd A\wedge\delta A)
  =
  \tr(\delta A\wedge\dd A).
  \label{eq:dense-var-commute}
\end{equation}
Putting these pieces together,
\begin{align}
  \delta\omega_{\CS}
  &=
  2\tr(\delta A\wedge\dd A)
  +
  2\tr(\delta A\wedge A\wedge A)
  -
  \dd\,\tr(A\wedge\delta A)
  \notag\\
  &=
  2\tr(\delta A\wedge F)
  -
  \dd\,\tr(A\wedge\delta A).
  \label{eq:dense-cs-variation-final}
\end{align}
Thus
\begin{equation}
  \delta S_{\CS}
  =
  \frac{k}{2\pi}\int_M\tr(\delta A\wedge F)
  -
  \frac{k}{4\pi}\int_{\partial M}\tr(A\wedge\delta A).
  \label{eq:dense-action-variation}
\end{equation}
The bulk equation is $F=0$, while the boundary term is the seed of the boundary-condition and anomaly-inflow story.

The boundary term should be kept in view rather than thrown away automatically. On a closed manifold it integrates to zero, but on a manifold with boundary it says that the bulk action alone does not define a complete variational problem. This is the first algebraic sign that the boundary theory in Section~\ref{sec:anomalies-boundaries-response} is not optional.

\subsection{Gauge Transformation}
\label{subsec:dense-gauge-transform-cs}

Let
\begin{equation}
  A^g=g^{-1}Ag+g^{-1}\dd g,
  \qquad
  \theta=g^{-1}\dd g,
  \qquad
  \widetilde A=g^{-1}Ag.
  \label{eq:dense-gauge-notation}
\end{equation}
Then $A^g=\widetilde A+\theta$. The Maurer--Cartan equation says
\begin{equation}
  \dd\theta=-\theta\wedge\theta.
  \label{eq:dense-MC}
\end{equation}
One may compute
\begin{equation}
  \dd A^g
  =
  \dd\widetilde A-\theta\wedge\theta.
  \label{eq:dense-dAg}
\end{equation}
This is a useful compressed identity: the terms involving $\theta\wedge\widetilde A$ and $\widetilde A\wedge\theta$ cancel between differentiating $\widetilde A$ and differentiating $\theta$.

Now expand the first part of the Chern--Simons form:
\begin{align}
  \tr(A^g\wedge\dd A^g)
  &=
  \tr\left((\widetilde A+\theta)\wedge(\dd\widetilde A-\theta\wedge\theta)\right)
  \notag\\
  &=
  \tr(\widetilde A\wedge\dd\widetilde A)
  -
  \tr(\widetilde A\wedge\theta\wedge\theta)
  +
  \tr(\theta\wedge\dd\widetilde A)
  -
  \tr(\theta^{\wedge 3}).
  \label{eq:dense-transform-quadratic}
\end{align}
Similarly, after taking the trace,
\begin{align}
  \tr((A^g)^{\wedge 3})
  &=
  \tr(\widetilde A^{\wedge 3})
  +
  3\tr(\widetilde A\wedge\widetilde A\wedge\theta)
  +
  3\tr(\widetilde A\wedge\theta\wedge\theta)
  +
  \tr(\theta^{\wedge 3}).
  \label{eq:dense-transform-cubic}
\end{align}
Combining \eqref{eq:dense-transform-quadratic} and \eqref{eq:dense-transform-cubic},
\begin{align}
  \omega_{\CS}(A^g)
  &=
  \tr\left(\widetilde A\wedge\dd\widetilde A+\frac23\widetilde A^{\wedge 3}\right)
  +
  2\tr(\widetilde A\wedge\widetilde A\wedge\theta)
  \notag\\
  &\quad
  +
  \tr(\widetilde A\wedge\theta\wedge\theta)
  +
  \tr(\theta\wedge\dd\widetilde A)
  -
  \frac13\tr(\theta^{\wedge 3}).
  \label{eq:dense-transform-combined}
\end{align}
The first term is not simply $\omega_{\CS}(A)$ because $\dd\widetilde A$ contains $\theta$ terms. A short substitution gives
\begin{equation}
  \tr\left(\widetilde A\wedge\dd\widetilde A+\frac23\widetilde A^{\wedge 3}\right)
  =
  \omega_{\CS}(A)
  -
  2\tr(\widetilde A\wedge\widetilde A\wedge\theta).
  \label{eq:dense-transform-tilde}
\end{equation}
The $\widetilde A^2\theta$ terms cancel. What remains is
\begin{equation}
  \omega_{\CS}(A^g)
  =
  \omega_{\CS}(A)
  +
  \tr(\widetilde A\wedge\theta\wedge\theta)
  +
  \tr(\theta\wedge\dd\widetilde A)
  -
  \frac13\tr(\theta^{\wedge 3}).
  \label{eq:dense-transform-almost}
\end{equation}
The middle two terms are exact:
\begin{equation}
  \tr(\widetilde A\wedge\theta\wedge\theta)
  +
  \tr(\theta\wedge\dd\widetilde A)
  =
  \dd\,\tr(\widetilde A\wedge\theta)
  =
  \dd\,\tr(g(\dd g^{-1})\wedge A).
  \label{eq:dense-transform-exact}
\end{equation}
Therefore
\begin{equation}
  \omega_{\CS}(A^g)
  =
  \omega_{\CS}(A)
  -
  \frac13\tr(\theta^{\wedge 3})
  +
  \dd\,\tr(g(\dd g^{-1})\wedge A).
  \label{eq:dense-transform-final}
\end{equation}
This formula is the precise place where local gauge covariance and global topology meet.

The exact term and the Wess--Zumino term should be mentally separated. The exact term is controlled by the boundary of spacetime. The cubic term is controlled by the homotopy class of the gauge transformation. That split is why the same transformation law feeds both anomaly inflow and level quantization.

\subsection{Winding Integral}
\label{subsec:dense-euler-su2}

For $SU(2)$, the generator of $\pi_3(SU(2))\cong\mathbb Z$ is the identity map $SU(2)\to SU(2)$. The normalization of the winding number is fixed by
\begin{equation}
  \frac{1}{24\pi^2}
  \int_{SU(2)}
  \tr\left((g^{-1}\dd g)^{\wedge 3}\right)
  =
  1.
  \label{eq:dense-su2-normalization-goal}
\end{equation}
Use Euler angles
\begin{equation}
  g(\psi,\theta,\phi)
  =
  e^{i\psi\sigma^3/2}
  e^{i\theta\sigma^2/2}
  e^{i\phi\sigma^3/2},
  \qquad
  0\leq\psi<4\pi,\quad
  0\leq\theta\leq\pi,\quad
  0\leq\phi<2\pi.
  \label{eq:dense-euler-param}
\end{equation}
A direct computation gives
\begin{align}
  g^{-1}\dd g
  &=
  \frac{i}{2}\sigma^1
  \left(\sin\theta\cos\phi\,\dd\psi+\sin\phi\,\dd\theta\right)
  \notag\\
  &\quad
  +
  \frac{i}{2}\sigma^2
  \left(-\sin\theta\sin\phi\,\dd\psi+\cos\phi\,\dd\theta\right)
  \notag\\
  &\quad
  +
  \frac{i}{2}\sigma^3
  \left(\cos\theta\,\dd\psi+\dd\phi\right).
  \label{eq:dense-ginv-dg-euler}
\end{align}
Define real one-forms
\begin{align}
  e^1&=\sin\theta\cos\phi\,\dd\psi+\sin\phi\,\dd\theta,
  \notag\\
  e^2&=-\sin\theta\sin\phi\,\dd\psi+\cos\phi\,\dd\theta,
  \notag\\
  e^3&=\cos\theta\,\dd\psi+\dd\phi.
  \label{eq:dense-left-invariant-forms}
\end{align}
Then
\begin{equation}
  g^{-1}\dd g=\frac{i}{2}\sigma^a e^a.
  \label{eq:dense-ginv-dg-compact}
\end{equation}
Using $\tr(\sigma^a\sigma^b\sigma^c)=2i\epsilon^{abc}$,
\begin{align}
  \tr\left((g^{-1}\dd g)^{\wedge 3}\right)
  &=
  \left(\frac{i}{2}\right)^3
  e^a\wedge e^b\wedge e^c
  \tr(\sigma^a\sigma^b\sigma^c)
  \notag\\
  &=
  \frac32\,e^1\wedge e^2\wedge e^3.
  \label{eq:dense-wz-e-forms}
\end{align}
The remaining wedge product is
\begin{align}
  e^1\wedge e^2
  &=
  \sin\theta\,\dd\psi\wedge\dd\theta,
  \notag\\
  e^1\wedge e^2\wedge e^3
  &=
  \sin\theta\,\dd\psi\wedge\dd\theta\wedge\dd\phi.
  \label{eq:dense-e123}
\end{align}
Thus
\begin{equation}
  \tr\left((g^{-1}\dd g)^{\wedge 3}\right)
  =
  \frac32\sin\theta
  \,\dd\psi\wedge\dd\theta\wedge\dd\phi.
  \label{eq:dense-wz-euler-density}
\end{equation}
Integrating over the Euler-angle ranges,
\begin{align}
  \int_{SU(2)}
  \tr\left((g^{-1}\dd g)^{\wedge 3}\right)
  &=
  \frac32
  \int_0^{4\pi}\dd\psi
  \int_0^\pi\sin\theta\,\dd\theta
  \int_0^{2\pi}\dd\phi
  \notag\\
  &=
  \frac32(4\pi)(2)(2\pi)
  =
  24\pi^2.
  \label{eq:dense-euler-integral}
\end{align}
This is the calculation that fixes the normalization in the level-quantization argument.

% FIGURE FLAG [F-su2-volume]: a figure of SU(2) as S^3 with Euler-angle ranges and the normalized volume form would clarify why the integral gives 24 pi^2.

\subsection{Torus Quantization}
\label{subsec:dense-u1-torus-quantization}

Let $\Sigma=T^2$ with coordinates $x^1,x^2\in[0,1)$ and take a flat connection
\begin{equation}
  A=u\,\dd x^1+v\,\dd x^2,
  \qquad
  u\sim u+2\pi,\quad v\sim v+2\pi.
  \label{eq:dense-flat-torus-A}
\end{equation}
The holonomies are $u=\oint_\alpha A$ and $v=\oint_\beta A$. Substituting into the symplectic form
\begin{equation}
  \Omega=\frac{k}{2\pi}\int_{T^2}\delta A\wedge\delta A
  \label{eq:dense-torus-symplectic-start}
\end{equation}
gives
\begin{equation}
  \Omega=\frac{k}{2\pi}\,\delta u\wedge\delta v.
  \label{eq:dense-torus-symplectic}
\end{equation}
The total symplectic volume is
\begin{equation}
  \int_{T^2_{\mathrm{hol}}}\Omega
  =
  \frac{k}{2\pi}
  \int_0^{2\pi}\dd u
  \int_0^{2\pi}\dd v
  =
  2\pi k.
  \label{eq:dense-torus-volume}
\end{equation}
Quantization counts one state per $2\pi$ unit of symplectic volume:
\begin{equation}
  \dim\mathcal H_{\U(1)_k}(T^2)
  =
  \int_{T^2_{\mathrm{hol}}}\frac{\Omega}{2\pi}
  =
  k.
  \label{eq:dense-torus-state-count}
\end{equation}

The conceptual step is the reduction from fields to holonomies. Before the constraint is imposed, $A_i(x)$ is a field on the whole spatial torus. After flatness and gauge equivalence, the only gauge-invariant coordinates left are the two periods around the fundamental cycles. Chern--Simons theory has converted an infinite-dimensional field problem into the quantization of a finite-dimensional phase space.

One can also see the same result from the Wilson-loop algebra. Let
\begin{equation}
  W_\alpha=e^{iu},
  \qquad
  W_\beta=e^{iv}.
  \label{eq:dense-wilson-holonomies}
\end{equation}
The symplectic form implies
\begin{equation}
  [u,v]=\frac{2\pi i}{k}.
  \label{eq:dense-u-v-commutator}
\end{equation}
Exponentiating,
\begin{equation}
  W_\alpha W_\beta
  =
  e^{2\pi i/k}W_\beta W_\alpha.
  \label{eq:dense-clock-shift}
\end{equation}
The smallest irreducible representation of this clock-shift algebra has $k$ states. If $\{|n\rangle\}_{n=0}^{k-1}$ is a basis, one may take
\begin{equation}
  W_\alpha|n\rangle=e^{2\pi i n/k}|n\rangle,
  \qquad
  W_\beta|n\rangle=|n+1\rangle,
  \label{eq:dense-clock-shift-rep}
\end{equation}
with $|k\rangle\equiv|0\rangle$. This is the operator version of the geometric-quantization count.

% FIGURE FLAG [F-torus-holonomies]: draw the alpha and beta cycles, the holonomy torus with coordinates (u,v), and the k-state clock-shift representation.

\subsection{Theta Wavefunctions}
\label{subsec:dense-theta-wavefunctions}

To make the $k$-dimensional Hilbert space more concrete, choose a complex structure $\tau$ on the holonomy torus and define
\begin{equation}
  z=\frac{1}{2\pi}(u+\tau v).
  \label{eq:dense-z-coordinate}
\end{equation}
The prequantum line bundle has curvature
\begin{equation}
  F_{\nabla}=-i\Omega=-i\frac{k}{2\pi}\dd u\wedge\dd v.
  \label{eq:dense-prequantum-curvature}
\end{equation}
Holomorphic wavefunctions are not ordinary periodic functions on the torus. They are sections of this line bundle, so under $u\mapsto u+2\pi$ and $v\mapsto v+2\pi$ they transform by phases determined by the connection. A basis can be written in terms of level-$k$ theta functions:
\begin{equation}
  \vartheta_\ell(z|\tau)
  =
  \sum_{n\in\mathbb Z}
  \exp\left[
  i\pi k\tau\left(n+\frac{\ell}{k}\right)^2
  +
  2\pi i k\left(n+\frac{\ell}{k}\right)z
  \right],
  \qquad
  \ell=0,\ldots,k-1.
  \label{eq:dense-theta-functions}
\end{equation}
The label $\ell$ is defined modulo $k$, so there are exactly $k$ independent sections. The theta-function construction is the wavefunction-level version of the same fact we saw from symplectic volume and from the Wilson-loop algebra.

These are wavefunctions on the moduli space of flat connections, not wavefunctions of a particle moving on the original spatial torus. That distinction is easy to miss. The theta functions know about the original torus only through the holonomies and the chosen complex structure on the reduced phase space.

\subsection{Wilson Loops}
\label{subsec:dense-abelian-wilson-gaussian}

For $\U(1)_k$ on a closed three-manifold, the Wilson loop of charge $q$ along $C$ is
\begin{equation}
  W_q(C)=\exp\left(iq\oint_C A\right).
  \label{eq:dense-abelian-wilson}
\end{equation}
Introduce the Poincare-dual two-form $\eta(C)$ by
\begin{equation}
  \int_C A=\int_M A\wedge\eta(C).
  \label{eq:dense-poincare-dual-loop}
\end{equation}
For several loops,
\begin{equation}
  J=\sum_\alpha q_\alpha \eta(C_\alpha).
  \label{eq:dense-source-J}
\end{equation}
The path integral becomes
\begin{equation}
  \left\langle\prod_\alpha W_{q_\alpha}(C_\alpha)\right\rangle
  =
  \frac{1}{Z}
  \int\mathcal DA\,
  \exp\left[
  i\frac{k}{4\pi}\int A\wedge\dd A
  +
  i\int A\wedge J
  \right].
  \label{eq:dense-gaussian-wilson}
\end{equation}
The saddle equation is
\begin{equation}
  \frac{k}{2\pi}\dd A+J=0.
  \label{eq:dense-gaussian-eom}
\end{equation}
Choose Seifert surfaces $\Sigma_\alpha$ with $\partial\Sigma_\alpha=C_\alpha$. Their Poincare duals obey
\begin{equation}
  \dd\eta(\Sigma_\alpha)=\eta(C_\alpha).
  \label{eq:dense-seifert-dual}
\end{equation}
A classical solution is therefore
\begin{equation}
  A_{\mathrm{cl}}
  =
  -\frac{2\pi}{k}
  \sum_\alpha q_\alpha\eta(\Sigma_\alpha).
  \label{eq:dense-A-classical}
\end{equation}
Substituting into the exponent gives
\begin{equation}
  \exp\left[
  \frac{\pi i}{k}
  \sum_{\alpha,\beta}
  q_\alpha q_\beta
  L_{\alpha\beta}
  \right],
  \label{eq:dense-linking-answer}
\end{equation}
where
\begin{equation}
  L_{\alpha\beta}
  =
  \int_M\eta(\Sigma_\alpha)\wedge\eta(C_\beta)
  \label{eq:dense-linking-intersection}
\end{equation}
is the intersection number of $\Sigma_\alpha$ with $C_\beta$, i.e. the linking number when $\alpha\neq\beta$. For two disjoint loops,
\begin{equation}
  \langle W_{q_1}(C_1)W_{q_2}(C_2)\rangle
  =
  \exp\left[
  \frac{2\pi iq_1q_2}{k}\Lk(C_1,C_2)
  \right].
  \label{eq:dense-two-loop-linking}
\end{equation}
The factor of $2$ arises because the off-diagonal entries $(1,2)$ and $(2,1)$ both appear in the symmetric sum in \eqref{eq:dense-linking-answer}.

The physical picture is an Aharonov--Bohm phase with no local propagating force mediating it. A Wilson loop inserts charge, the Chern--Simons field attaches flux to that charge, and another loop detects the flux through its linking number. Because the pure theory has no metric-dependent bulk dynamics, the answer depends on topology rather than on the detailed shapes of the curves.

% FIGURE FLAG [F-wilson-linking]: show two linked loops, a Seifert surface for one loop, and its intersection with the other loop.

\subsection{Framing}
\label{subsec:dense-framing}

The diagonal term $L_{\alpha\alpha}$ is not defined by an ordinary linking number of two disjoint curves, because there is only one curve. To define it, choose a framing: a nonzero normal vector field along $C$. Pushing $C$ slightly in that normal direction gives a nearby curve $C'$, and the self-linking number is
\begin{equation}
  \SLk(C)=\Lk(C,C').
  \label{eq:dense-self-linking}
\end{equation}
Changing the framing by one twist changes $\SLk(C)$ by one. In the abelian Wilson-loop answer, this changes the expectation value by
\begin{equation}
  W_q(C)
  \mapsto
  \exp\left(\frac{\pi i q^2}{k}\right)W_q(C)
  \label{eq:dense-framing-shift}
\end{equation}
for a unit change in self-linking, with the precise sign depending on orientation conventions. This is why Chern--Simons theory naturally computes framed link invariants. Physically, the same phase is interpreted as topological spin data.

% FIGURE FLAG [F-framing-ribbon]: show a framed knot as a ribbon and the push-off curve whose linking number defines self-linking.

\subsection{\texorpdfstring{$K$}{K}-Matrix Theory}
\label{subsec:dense-kmatrix-integrateout}

The general abelian action is
\begin{equation}
  S_K[a^I;A]
  =
  \int\left[
  \frac{1}{4\pi}K_{IJ}a^I\wedge\dd a^J
  +
  \frac{1}{2\pi}t_I A\wedge\dd a^I
  \right].
  \label{eq:dense-K-action}
\end{equation}
Varying with respect to $a^I$ gives
\begin{equation}
  \frac{1}{2\pi}K_{IJ}\dd a^J
  +
  \frac{1}{2\pi}t_I\dd A
  =
  0.
  \label{eq:dense-K-eom}
\end{equation}
Thus
\begin{equation}
  \dd a^I
  =
  -(K^{-1})^{IJ}t_J\,\dd A.
  \label{eq:dense-K-solution}
\end{equation}
Substituting back, the electromagnetic response action is
\begin{equation}
  S_{\mathrm{eff}}[A]
  =
  -\frac{1}{4\pi}
  t^{\mathsf T}K^{-1}t
  \int A\wedge\dd A,
  \label{eq:dense-K-effective}
\end{equation}
up to the sign convention for the original $K$-matrix action. Therefore
\begin{equation}
  \sigma_{xy}
  =
  \frac{e^2}{h}\,t^{\mathsf T}K^{-1}t.
  \label{eq:dense-K-hall}
\end{equation}
For a quasiparticle labeled by an integer vector $\ell$, coupling $\ell_I a^I$ to a worldline gives charge
\begin{equation}
  Q_\ell=e\,t^{\mathsf T}K^{-1}\ell
  \label{eq:dense-K-charge}
\end{equation}
and mutual braiding phase
\begin{equation}
  \theta_{\ell\ell'}
  =
  2\pi\,\ell^{\mathsf T}K^{-1}\ell'.
  \label{eq:dense-K-mutual}
\end{equation}
The torus degeneracy is
\begin{equation}
  \dim\mathcal H(T^2)=|\det K|,
  \label{eq:dense-K-torus-degeneracy}
\end{equation}
and on genus $g$ it becomes
\begin{equation}
  \dim\mathcal H(\Sigma_g)=|\det K|^g.
  \label{eq:dense-K-genus}
\end{equation}
This single matrix package is the sense in which Chern--Simons theory is a framework rather than one model.

The matrix inverse appears everywhere because integrating out the internal gauge fields solves a linear response problem. The matrix $K$ pairs emergent fluxes with emergent charges; $K^{-1}$ tells us how much internal flux is induced by an electromagnetic field or by a quasiparticle source. That is why the same algebra controls Hall conductance, fractional charge, braiding, and degeneracy.

% FIGURE FLAG [F-kmatrix-flow]: a compact flow diagram from K and t to Hall response, charge, mutual statistics, and ground-state degeneracy would connect the formulas.

\subsection{Chiral Edge Boson}
\label{subsec:dense-boundary-chiral-boson}

On a manifold with boundary, the flatness equation in the abelian theory says locally
\begin{equation}
  A_i=\partial_i\phi.
  \label{eq:dense-boundary-pure-gauge}
\end{equation}
Substituting a pure-gauge spatial connection back into the Chern--Simons action leaves a boundary term. In coordinates $(t,x)$ on the boundary, one obtains the chiral kinetic structure
\begin{equation}
  S_{\partial}
  =
  \frac{k}{4\pi}\int_{\partial M}
  \partial_t\phi\,\partial_x\phi\,\dd t\,\dd x
  +\cdots.
  \label{eq:dense-boundary-kinetic}
\end{equation}
The dots include nonuniversal Hamiltonian terms such as
\begin{equation}
  -\frac{k v}{4\pi}\int_{\partial M}
  (\partial_x\phi)^2\,\dd t\,\dd x,
  \label{eq:dense-boundary-hamiltonian}
\end{equation}
where $v$ is an edge velocity. The topological bulk fixes the anomaly coefficient and chirality; it does not fix the microscopic velocity.

Under a background gauge transformation, the chiral boson action shifts by the negative of the bulk boundary variation. This is the operational meaning of anomaly inflow in the quantum Hall setting. The edge mode is required not because the bulk has local degrees of freedom, but because gauge invariance of the full bulk-plus-boundary system demands it.

% FIGURE FLAG [F-inflow-edge]: a strip or disk geometry showing bulk Chern--Simons variation canceled by a chiral boundary mode would be useful in the anomaly section.

\subsection{Topological Mass}
\label{subsec:dense-mcs-mode}

Consider
\begin{equation}
  S_{\mathrm{MCS}}
  =
  \int\dd^3x
  \left[
  -\frac{1}{4g^2}f_{\mu\nu}f^{\mu\nu}
  +
  \frac{k}{4\pi}\epsilon^{\mu\nu\rho}a_\mu\partial_\nu a_\rho
  \right].
  \label{eq:dense-mcs}
\end{equation}
Varying with respect to $a_\mu$ gives
\begin{equation}
  \partial_\nu f^{\nu\mu}
  +
  \frac{k g^2}{4\pi}
  \epsilon^{\mu\nu\rho}f_{\nu\rho}
  =
  0.
  \label{eq:dense-mcs-eom}
\end{equation}
Define the dual field strength
\begin{equation}
  F^\mu=\frac12\epsilon^{\mu\nu\rho}f_{\nu\rho}.
  \label{eq:dense-mcs-dual}
\end{equation}
Then $f_{\nu\rho}=\epsilon_{\nu\rho\sigma}F^\sigma$, and the equation becomes
\begin{equation}
  \epsilon^{\nu\mu\sigma}\partial_\nu F_\sigma
  +
  \frac{k g^2}{2\pi}F^\mu
  =
  0.
  \label{eq:dense-mcs-first-order}
\end{equation}
Acting once more with $\epsilon_{\alpha\beta\mu}\partial^\beta$ and using the Bianchi identity $\partial_\mu F^\mu=0$, one obtains
\begin{equation}
  (\Box+m^2)F_\alpha=0,
  \qquad
  m=\frac{|k|g^2}{2\pi}.
  \label{eq:dense-mcs-mass}
\end{equation}
Thus the Chern--Simons term, when paired with a metric-dependent Maxwell term, produces a genuine propagating massive mode. Pure Chern--Simons theory is topological; Maxwell--Chern--Simons theory is not. The contrast sharpens the meaning of both.

The lesson is not simply that Chern--Simons terms can make particles massive. The lesson is that the Maxwell term supplies the local degree of freedom and the metric, while the Chern--Simons term changes the first-order structure of its equations of motion. Removing the Maxwell term removes the propagating mode, leaving only the topological sector.
